Process-based drainage and water-quality models---DRAINMOD and its
derivatives, SWAT, APEX, RZWQM2, HYDRUS, MIKE~SHE, and numerous salinity
codes---are expected to carry much of the predictive burden as field trials
cannot span every combination of practice, soil, and climate. Yet across the
literature their credibility is constrained less by their formulation than by
the narrowness of the conditions under which they have been tested: models are
typically calibrated at a single site, over a handful of seasons, for one
soil--climate--management combination, and their authors are nearly unanimous
in flagging broader empirical validation, missing process data, and weak
algorithmic components as the outstanding priorities
\parencite{Li2006,Qi2015,Youssef2006,Radcliffe2015,SKAGGS1994}. This section
gathers the model-algorithm and model-structure gaps that recur throughout the
field.

\subsection{Broadening the empirical basis for validation}

The most frequently articulated need is validation data spanning a wider range
of conditions than any single study can supply. \textcite{Youssef2006} describe
rigorous field testing across differing management, soils, and climates as a
need common to nearly all simulation models, and warn that assumptions must be
re-verified whenever a code is transferred to a new area
\parencite{Kelleners2000}; similar calls for broader-condition testing,
sensitivity analysis, and parameter refinement recur for pesticide, salt, and
flow models \parencite{Branger2009,Hirekhan2007,Yang2008b,Xian2017}. Models are
repeatedly reported as requiring testing over longer periods and additional soil
types \parencite{Salazar2009,Li2026}, and their effectiveness under new
controlled-drainage or erosion conditions is judged critically needed before
higher credibility can be claimed \parencite{Negm2017,Turunen2020}.

A recurring, more ambitious hypothesis is that empirical coefficients could be
replaced with functions of measurable site attributes. \textcite{Li2006} commit
to testing across climate zones, soils, and management while anticipating that
coefficient constants might become functions of soil texture, converting
site-specific tuning into transferable structure; \textcite{Chelil2022} cast the
same concern as spatial robustness, arguing that additional fields are needed to
confirm that parameters such as nitrate-transfer velocities are transferable
rather than calibration artefacts. Verification at independent sites is treated
as a precondition, not optional confirmation
\parencite{Azevedo1997,Qi2015,Kalita1998,Sadhukhan2019b}, with inconsistencies
producing conflicting parameter values resolved first \parencite{Dar2022} and
empirical codes transferred only through re-parameterisation with local soil,
geology, crop, and climate data, or region-specific coefficients where reference
values exist \parencite{Azevedo1997,FathAllah2025,Rijal2013}. The open question
is how much local data such transfer requires. Extending validation across the
climatic and design envelope is closely related: humid-region models need
testing under semi-arid regimes and after prolonged dry spells
\parencite{Prasher1996}, longer monitoring records give firmer conclusions
\parencite{Ebrahimian2011,Grazhdani1996}, and design geometry demands
experiments at spacings wider than 30~m and regional answers on
drainage-efficiency loss and design storms
\parencite{Mathew2001,Wiskow2003,MULLER1988}. Long simulations of 50--100 years
are argued necessary to judge whether a design meets its criteria
\parencite{Mendez2004,Wang2006b,Ale2009,Häggblom2019}. Portability, in short, is
a hypothesis to be tested, not a property to be presumed
\parencite{Hornbuckle2012,Shipitalo2016,Yousfi2014}.

\subsection{Field and process data still needed to constrain models}

Beneath the demand for broader validation lies a scarcity of process-level
measurements. Nitrogen dynamics emerge as the least well-constrained component:
\textcite{David2009} argue that systematic denitrification measurements are
required to test and refine models, their ideal being a multi-year,
tile-drained corn--soybean study measuring every major N flux, with a pragmatic
fallback of a few key fluxes plus denitrification. Field- and
laboratory-grounded parameterisation, denitrification foremost, is likewise
judged necessary \parencite{Salazar2009}, together with routine evaluation and
reporting of soil-N fluxes in SWAT \parencite{Ikenberry2017}, long-term
carbon--nitrogen estimates to reconcile plot- and subbasin-scale discrepancies
\parencite{Singer2011}, and resolution of the lag between practice
implementation and measurable reduction \parencite{Law2022,Tao2021}.
High-frequency process measurements of microbial activity, denitrification,
soil-P pools, and redox are urged for drainage water management, alongside
closure of water and nutrient balances \parencite{Frankenberger2024,Sharpley2019}.

Phosphorus is similarly data-limited, and its dominant pathway---preferential
flow---is the hardest to observe. Many phosphorus-index factors need
investigation of their effect on field-to-stream transfer, and index rankings
must be shown to reflect actual transfer risk
\parencite{Bechmann2007,Beauchemin2003}. Confirming P-movement trends requires
quantitative tracer measurement of preferential flow
\parencite{Salazar2011,Tiemeyer2009}, structural deficiencies around surface
runoff and tile P transport need repair \parencite{Apostel2021}, and accurate
capture of flashy export demands automated samplers, real-time sensors, and
high-resolution topographic survey of ditch sediment before and after storms
\parencite{Dialameh2022,Warsta2014}. Supporting empirical work includes further
testing of CaCl$_2$-P as a dissolved-reactive-P estimator \parencite{McDowell2001b},
regression models linking drainage parameters to P losses \parencite{Morrison2013},
tests using actual effluent under fluctuating flow \parencite{Sanford2016}, and
long-term monitoring of P sources \parencite{Jia2012}. Sampling design is itself
a constraint: weekly sampling may be needed to capture nitrate spikes before
they pass below suction-cup samplers \parencite{Zotarelli2007}, more frequent
early-season and summer-event sampling is advised
\parencite{Kalita1998,Tiemeyer2009}, and denser sampling geometry sharpens
semi-variograms and optimal sampling density \parencite{GALLICHAND1992,Upadhyay2019}.

\subsection{Preferential and macropore flow representation}

The representation of preferential flow is perhaps the most frequently cited
frontier, the recurring theme being that its contribution to drains must be
characterised as a function of soil type, water content, and practice before
models extrapolate. \textcite{Gaillot2023} and \textcite{Fox2004} call for
mechanistic characterisation of macropore--drain connectivity and of the
``express fraction'' under varying saturation above the drain, and
\textcite{Askar2020} stress that rigorous testing of DRAINMOD's empirical
macropore approach requires water-table-depth measurements alongside drainage
and runoff, direct macropore-flow inputs, and crack geometry and dynamics.
Process improvements span better matrix--macropore partitioning calibrated
against sub-hourly drainage in RZWQM2 \parencite{Li2023d}, dynamic connectivity
or effective macroporosity at watershed scale \parencite{Poon2021}, micropore
diffusion and site-specific conductivities in dual-permeability descriptions
\parencite{Villholth1998,Villholth1998b}, hysteresis and biopore-flow initiation
\parencite{Holbak2021}, manure-driven macropore development \parencite{Zhao2001},
crack re-establishment in EPIC \parencite{Wang2022}, and calibrated macropore
conductivity in three-dimensional dual-permeability models \parencite{Varvaris2018}.
Direct pathway identification is repeatedly urged---ground-penetrating radar and
tracers with explicit flowpath representation \parencite{Tähtikarhu2025,Salo2017},
dye-staining of routes reaching tile lines \parencite{Nielsen2015}, and
near-surface pathway characterisation \parencite{Frey2012,Allaire2011}. Central
quantitative questions remain the fraction of phosphorus and solutes delivered
via preferential versus matrix flow \parencite{King2015,Williams2016}, the range
of macropore flows across soils and management \parencite{Reid2012}, and by-pass
flow through cracks for full-year application
\parencite{LESAFFRE1988,BACKLUND1995,Ruark2009}. Runoff simulation under
crusting, cracking, and preferential flow \parencite{Smith2020}, explicit
representation of actual drain configuration \parencite{Stamm2002}, temporal
variability in hydrological-pathway activation for cracking clays
\parencite{Kokulan2019}, and seasonal, time-dependent shrink--swell and
hydraulic-conductivity measurement \parencite{Boivin2006,Chelil2022b,Jeantet2021}
complete this agenda.

\subsection{Coupled hydraulics, seepage, and soil heterogeneity}

A recurrent structural gap is the absence of models coupling soil pore-water
flow with surface flow inside drains and ditches. \textcite{Xin2016b} call for
coupling drain surface water with pore-water flow to represent the time-varying
head around the drain, and \textcite{Nousiainen2015} propose simulating in-pipe
flow to capture pressurisation and clogging, adding open-ditch networks.
Pressurised-flow dynamics are a specific frontier: coupling pipes to the water
table could introduce a validated temporary-retention function beyond design
capacity \parencite{Henine2010,Henine2014}, with pressurised periods analysed
separately from pressure-transducer data \parencite{Saadat2018}. Entrance
resistance, near-drain pressure heads, and rapid trench-backfill versus slower
inter-trench transport remain unresolved for heterogeneous systems
\parencite{Messing2006,Rimidis2003}, and model dimensionality is itself a
question---two- or three-dimensional treatment of sloping fields, collectors,
and regional management \parencite{Alakukku2010,FIPPS1989,Carlier2007,Mohanty1998}---within
broader integrating frameworks \parencite{Manguerra1998,SAVABI1993,RevueltaAcosta2021}.

Nearly every modelling study attributes residual error to unresolved
heterogeneity in hydraulic properties, especially conductivity. Authors seek
detailed representation of macropores, fractures, and sand lenses with specific
storage and residual saturation \parencite{DeSchepper2017}, depth-varying and
dual-porosity parameterisations in layered and paddy soils
\parencite{Wesström2003,Karandish2018}, and denser, depth-resolved
sampling---at least one conductivity measurement per hectare with
post-construction verification \parencite{PalaciosVélez2004}, more transect
points \parencite{Messing2006}, and horizontal as well as vertical
conductivities with hydrophobicity and hysteresis \parencite{Tähtikarhu2025}.
Anisotropy and long-term property change form a related sub-frontier
\parencite{Akis2015,Moustafa1998,Koivusalo1999}, and accurate saturated
conductivity and drainage-volume/upward-flux relationships, supported by
SSURGO-linked subsoil porosity, are emphasised as model inputs
\parencite{DeBoer2006,Schilling2015,DeSchepper2015}. Closing the water
balance---particularly seepage and groundwater exchange---remains persistent:
lateral seepage should vary during simulation \parencite{Feng2018}, groundwater
contributions could enter through double- or residual-storage models
\parencite{DaneshYazdi2016}, and groundwater-table data in objective functions
would sharpen inverse recovery of soil properties
\parencite{Sedaghatdoost2018,Sedaghatdoost2019}. Improved calibration strategy is
widely urged---parameter-estimation software and multiple conceptual models
\parencite{Boico2022}, pilot-point conductivity fields with proxy
machine-learning predictions \parencite{Schneider2022}, Kling--Gupta objective
functions \parencite{Brighenti2023}, and transient tests of steady-state
assumptions \parencite{Schilling2022,Gates2002,Burkhalter2006}. Contemporary
needs also span unsaturated modules for disconnected groundwater
\parencite{Kompanizare2024}, independent layer-to-layer drainage methods under
preferential flow \parencite{Ojeda2018}, and dynamic frozen-soil models with
measured structure \parencite{Ulén2016}.

\subsection{Nutrient routines, cold-season, and groundwater coupling}

Phosphorus and its transport through macropores emerge as the least mature
component across water-quality models. \textcite{Radcliffe2015} provide the most
comprehensive agenda: nested paired field experiments and structural additions
across codes---SWAT and ADAPT routing macropore flow directly to drains, a
macropore component for APEX, DRAINMOD simulating P losses and sediment-attached
P, and HYDRUS parameterising simpler models---with every study quantifying
predicted-P accuracy and uncertainty. Consistent with this,
\textcite{Guo2020b} recommend physically based tile subroutines, macropore
nutrient pathways, and particulate-P losses, \textcite{Ford2017} seek
understanding of how macropore flow and P respond to tillage, cover crops, and
drainage water management, and \textcite{BinShokrana2023b} identify improved
fertiliser-P partitioning and P-only performance statistics, with
\textcite{Gowda2012} anticipating a DRAINMOD phosphorus module. Within SWAT,
APEX, and RZWQM2, \textcite{Williams2016} argue that incorporating identified
flow pathways is critical, \textcite{Qi2018} recommend a soil-matrix-potential
Richards equation and revised particulate-P algorithm, and
\textcite{Moriasi2012} call for comparing dynamic and constant surface-depression
storage. Nitrogen-cycle needs include surface-residue and mobile/immobile-water
phases for nitrate transport \parencite{Malone2017}, complete N budgets in
calibration \parencite{Gillette2018}, and temporal decay coefficients
\parencite{Rudra2005}.

Soil N-process representation is repeatedly the weakest link. Priorities include
improving nitrification, denitrification, and mineralisation within a better
hydrology framework \parencite{Smith2020,Fang2012,Hu2007}, and recurring RZWQM
agendas---subsurface lateral flow below tiles, simulated soybean N fixation,
improved plant N uptake, and reduced winter nitrification
\parencite{Ma2007,Ma2008,Ma2001,Bakhsh2004}---together with revised winter
processes in the PNM model \parencite{Marjerison2016} and a more physically based
tile-nitrate model using ``soft data'' and a passive humus pool
\parencite{Qi2025}. Cold-season representation is singled out as inadequate:
better freeze/thaw and spring-thaw algorithms \parencite{Feyereisen2006b},
DRAINMOD-N improvement for freezing seasons \parencite{Yang2007,Saadat2020,Jiang2020},
a frozen-soil impedance factor and calibratable snow--rain threshold for RZ-SHAW
\parencite{Li2024b}, and continued RZWQM winter development
\parencite{Ahmed2007b,Motarjemi2021}. Coupling drainage models to groundwater is
a distinct need: SWAT-MODFLOW-RT3D for reactive transport once salinity is
incorporated \parencite{Addab2022}, water-table-fluctuation and
groundwater--surface-water exchange in SWAT to represent riparian nitrate
retention \parencite{Hoang2017}, joint spatial variation in leaching and
reduction across systems \parencite{Hojberg2017}, and defensible water balances
splitting artificial drainage from groundwater flow
\parencite{Barkle2021,Elsayed2023,Yang2008}. In arid settings, coupling of water
and salt dynamics with unsaturated models, root-growth modules, and
large-dataset validation is pressed
\parencite{Ebrahimian2015,Sarangi2006,Han2025,Liu2022,Guo2025}.

\subsection{Calibration, sensitivity analysis, and uncertainty}

How models are calibrated and their parameters interrogated is inseparable from
validity. Multiobjective frameworks integrating statistical and graphical
diagnostics tailored to the hydrological system are advocated
\parencite{Bruhn2025}, reliable soil-hydraulic data for the drain layer is
emphasised with interacting parameters calibrated and negligible ones fixed
\parencite{Wang2006}, and independent calibration against observed discharge is
recommended so interception and storage are not conflated with conductivity and
drainable porosity \parencite{Henine2022}. Pedotransfer estimates such as
Rosetta conductivity should be used cautiously, with further calibration so both
water-table depth and discharge are simulated \parencite{Malota2022,Timalsina2023}.
A common thread is identifiability---separating parameters that genuinely act
from those merely correlated. Automated and global approaches are widely urged:
automatic calibration with ensemble Kalman filtering for parameter identification
\parencite{Warsta2014,Sadhukhan2019b,Singh2019b}, global sensitivity analysis to
capture interactions that local one-at-a-time methods miss
\parencite{Timalsina2023,Chelil2022,Kacimov2000}, established bounds for the
SWAT sensitivity adjustment factor \parencite{Moriasi2013b}, and spatially
distributed streambed conductance with exchange fluxes as calibration targets and
drainage proxies \parencite{Karan2021}. Explicit uncertainty treatment is a
distinct, underdeveloped frontier: propagation of parameter and measurement error
\parencite{Alzraiee2013b}, the role of input-data and weather-data quality
\parencite{Upadhyay2019,Tao2021}, explicit estimation of P-model prediction
accuracy \parencite{Radcliffe2015,Beauchemin2003}, structural uncertainty via
multiple water-quality models \parencite{Kujawa2020}, probabilistic
drainage-water-management effectiveness \parencite{Shedekar2021}, independent
evaluation against withheld data \parencite{Schoups2005}, and better statistical
treatment of unreplicated paired experiments \parencite{Pagliari2021,Andrus2014}.
The recurring message is that models are too often reported as point predictions
when their inputs warrant propagated bounds.

\subsection{Model structure, coupling, and machine learning}

Beyond recalibration, structural refinements and integrative strategies are
judged necessary. Coupling complementary models is a common, largely untested
proposal: linking biogeochemical to spatial hydrological models such as SWAT or
MIKE~SHE for regional prediction \parencite{Li2006}, coupling nitrate-transport
and drainage models with variational assimilation of spring-fertilisation
information \parencite{Chelil2022}, and integrated use of more than one model as
more practical than a single code \parencite{Singh2019b,Apostel2021}.
Machine-learning models raise their own validation requirements the literature
is only beginning to confront: hyperparameter tuning and diverse training data
for drain detection \parencite{Song2021}, neural networks needing training that
includes extreme situations and periodic retraining \parencite{Yang2000}, and
machine-learning surrogates trained across diverse climatic and hydrogeological
conditions with transient processes and real-time control
\parencite{Lei2026,Schneider2022}. Process-model deficiencies are named
directly---detailed site characterisation, added non-linear regression
dependences, and improved submodels for mineralisation, crop N uptake, and
growth under N limitation \parencite{Ahuja2007,Ikenberry2017}, with added
complexity the price of predicting snowmelt-thaw drainage timing
\parencite{Feyereisen2007}. Physical and analytical refinements form a rich set
of concrete proposals: numerical models incorporating anisotropy and varying
bottom slopes verified against analytical solutions \parencite{Lin2025,Kacimov2000},
reanalysis with the Laplace rather than Boussinesq equation \parencite{Bouarfa2000},
recession equations extended to spatially varying properties \parencite{Cooke2001},
three-dimensional treatment at heterogeneous sites \parencite{Mohanty1997},
distinct mixing-layer depths and partition coefficients for solute transport
\parencite{Shalit1996}, reconsidered deep-percolation and shallow-water-table
assumptions \parencite{Kolars2019}, and strengthened index-system construction
with additional soil indicators \parencite{Hou2018}.

\subsection{Instrumentation and measurement-method validation}

Because model validation is only as sound as the data feeding it, uncalibrated
flow structures and drifting sensors can masquerade as model inadequacy.
In-house calibration of flow-measurement structures is recommended repeatedly:
weir calibration for individual applications \parencite{Christianson2019,Napierala2023},
in-house calibrations representing field flow conditions for control structures
\parencite{Katuwal2024}, stage--discharge equations developed before monitoring
for cylindrical structures \parencite{BinShokrana2021}, and post-hoc sensor
calibration against grab samples during tracer tests \parencite{deOliveira2023b}.
Emerging technologies carry explicit validation requirements: a computer-vision
flow method needs testing under variable real-world conditions with the
Manning--Strickler equation to compensate for uneven velocities
\parencite{Young2023}, hybrid RTK--rotating-laser systems address
less-than-ideal survey conditions \parencite{Freeland2014}, and sensor design
and scintillometer humidity-stability effects require refinement
\parencite{Enamorado2007,Zhang2010}. Hydraulic-loss characterisation in drainage
conduits remains conspicuously incomplete, limiting design-equation realism:
head loss in corrugated-pipe junctions and manholes under free-surface flow
\parencite{Bayat2011}, friction loss and transitional cavity-flow conditions
\parencite{Nédélec2008}, declining-bend slope effects on air discharge
\parencite{Gurklys2008}, appropriate Manning coefficients for design
\parencite{Kouchakzadeh2011}, and strain-gauge characterisation of non-linear
HDPE-pipe behaviour \parencite{Gaj2021}. The consistent message across the
modelling frontier is that the next decade of progress depends less on new
equations than on sustained, well-instrumented, and openly reported field
campaigns needed to test the equations already in hand.